\section{Methodology}

\subsection{Overview}

Explaintool processes a legal document through five layers: ingestion,
knowledge-graph construction, embedding and indexing, retrieval and
explanation, and answer generation. The retrieval layer (HHGR) is the
methodological core of this paper.

\subsection{Hierarchical knowledge graph}

Each statute is parsed into a hierarchy of documents, chapters, sections,
explanations, illustrations, and provisos. These become nodes in a knowledge
graph with \texttt{PART\_OF} edges encoding the document structure, plus
\texttt{CITES} and \texttt{REFERENCES} edges for intra- and inter-document
citations. For every node we store its title, text, numbering, and hierarchy
level, enabling structural reasoning that flat chunking cannot express.

\subsection{Dense multilingual retrieval}

We index every node with a multilingual embedding model. Retrieval is performed
by cosine similarity in the shared vector space, so queries written in Hindi
(for example \emph{अनुबंध प्रदर्शन}, ``contract performance'') can match English
evidence without explicit translation. In the released offline package the
provider is a deterministic hashing embedder so results are fully reproducible
without a model download; production deployments swap in a real multilingual
model such as bge-m3.

\subsection{Graph retrieval}

HHGR's graph path combines four signals per node: text overlap with the query,
citation matching against referenced section numbers, hierarchy position
(structural importance), and evidence propagated from dense matches to
ancestors and descendants along \texttt{PART\_OF} edges. This lets a query about
``performance'' surface both Section~4 and its Chapter~III ancestor, giving the
answer its statutory context.

\subsection{Fusion and explainability}

The final score of each candidate node is a weighted combination of the dense,
graph, and hierarchy signals, with default weights $\mathbf{w} = (0.40, 0.35,
0.25)$ normalized to sum to one. From the ranked evidence we construct:

\begin{itemize}
  \item \textbf{Evidence}: per-signal scores, active signal sources, and the
        root-to-node hierarchy path.
  \item \textbf{Reasoning chain}: a six-step audit trail (query parsing, dense
        search, graph search, hierarchy propagation, fusion, verification).
  \item \textbf{Citations}: numbered source citations with snippets.
  \item \textbf{Counter-authority detection}: marker-based detection of
        overruled, superseded, repealed, and inoperative authorities.
  \item \textbf{Confidence and validity}: a composite confidence score and
        boolean validity flags that gate the generated answer.
\end{itemize}

\subsection{Evaluation metrics}

Retrieval quality is measured with Recall@K, Precision@K, Hit Rate@K, MRR,
MAP, and NDCG@K, together with latency statistics (mean, p50, p95) and
throughput. Explainability is measured with citation accuracy, hierarchy
correctness, graph-path accuracy, provenance completeness, evidence coverage,
and counter-authority detection F1. Generation quality is measured with
RAGAS-style \cite{es2024ragas} faithfulness, answer relevancy, context recall,
context precision, and answer correctness, computed with deterministic offline
surrogates so the package runs without external services.
